\documentclass[11pt,a4paper]{article}
\usepackage[margin=27mm]{geometry}
\usepackage{amsmath,amsthm,amssymb,fontspec,unicode-math,hyperref}
\setmainfont{STIXGeneral}
\setmathfont{STIX Two Math}
\setmonofont{Menlo}
\hypersetup{hidelinks}
\usepackage{newunicodechar}
\newfontfamily\numeralfallback{STIX Two Text}
\newunicodechar{∎}{\ensuremath{\blacksquare}}
\newunicodechar{⁰}{{\numeralfallback\char"2070\relax}}
\newunicodechar{₀}{{\numeralfallback\char"2080\relax}}
\newunicodechar{₁}{{\numeralfallback\char"2081\relax}}
\newunicodechar{₂}{{\numeralfallback\char"2082\relax}}
\newunicodechar{₃}{{\numeralfallback\char"2083\relax}}
\newunicodechar{₄}{{\numeralfallback\char"2084\relax}}
\newunicodechar{⁻}{{\numeralfallback\char"207B\relax}}
\newunicodechar{⁸}{{\numeralfallback\char"2078\relax}}
\newunicodechar{ⁱ}{{\numeralfallback\char"2071\relax}}
\newunicodechar{⁴}{{\numeralfallback\char"2074\relax}}
\newunicodechar{⁷}{{\numeralfallback\char"2077\relax}}
\newunicodechar{⁶}{{\numeralfallback\char"2076\relax}}
\newunicodechar{⁵}{{\numeralfallback\char"2075\relax}}
\newunicodechar{₈}{{\numeralfallback\char"2088\relax}}
\newunicodechar{ᵢ}{\textsubscript{i}}
\newunicodechar{ⱼ}{\textsubscript{j}}
\newunicodechar{ₜ}{\textsubscript{t}}
\providecommand{\tightlist}{\setlength{\itemsep}{0pt}\setlength{\parskip}{0pt}}
\newtheorem{theorem}{Theorem}
\newtheorem{proposition}[theorem]{Proposition}
\newtheorem{lemma}[theorem]{Lemma}
\newtheorem{corollary}[theorem]{Corollary}
\theoremstyle{remark}\newtheorem{remark}[theorem]{Remark}
\newcommand{\gen}[1]{\langle #1\rangle}
\newcommand{\normalgen}[1]{\langle\!\langle #1\rangle\!\rangle}
\newcommand{\Z}{\mathbb Z}
\newcommand{\F}{\mathbb F}
\setlength{\emergencystretch}{2em}

\usepackage{longtable,booktabs,array,calc}
\newcounter{none}
\setlength{\emergencystretch}{4em}
\providecommand{\passthrough}[1]{#1}
\title{Cyclic cores and high-type two-group polyhedra}
\author{Human Compression Project}
\date{}
\begin{document}
\maketitle
\begin{abstract}
We recover the all-parameter existence theorem and the complete eight-family
high-type classification of Hou, Feng and Leemans. Existence follows by
attaching two dihedral quotients to a common affine square group. For
classification, a cyclic-core bound reduces every candidate to one of
three bounded square quotients. Explicit extension models give all
orders and markings; converse arguments and sharp presentation bounds
prove completeness. A single exhaustive order-256 kernel computation
supplies the remaining finite seed, without a group catalogue.
\end{abstract}
\section{The source results and common inputs}
We recover Theorems 1.2 and 1.3 of Hou, Feng and Leemans
\cite{HFL}, in the author edition \nolinkurl{arXiv:1802.08947v2} dated
22 January 2019. The first theorem is an existence statement for every
admissible rank-three type. The second classifies three high-type
families, with two, two and four marked groups respectively. A proof of
existence alone does not establish this classification.

We use the marking conventions of the companion common paper
\cite{UA}. In particular, a string C-group is generated by nonidentity
involutions, distant generators commute, and subgroups on two subsets
of the generators intersect in the subgroup on their intersection.
For rank three we write the marking as $(a,b,c)$, with $ac=ca$.
Conjugation is on the right: $x^g=g^{-1}xg$, and
$[x,g]=x^{-1}g^{-1}xg$. The subscript of $D_{2m}$ denotes its order.
All isomorphisms used to classify the models preserve the ordered
marking. The standard correspondence with regular abstract polytopes
is retained from \cite[Section 2E]{MS}.

The reusable results used below are UA.Phi, UA.Square,
UA.SquareRecognition, UA.Edge and UA.Search, the last including the
complete UA.Nilpotent and UA.Ladder proofs. They are fully stated and
proved in \cite{UA}. The application-specific work consists of the
cyclic-core reduction, eight explicit extension models, their converse
classification, and recognition of the bare source presentations. The
only numerical input is a complete order-256 marked-kernel computation,
whose generation is justified by UA.Search. Its full program is included
in the shared computational appendix; no abstract-group catalogue or
historical classification count is an input.

\section{Existence for every stated type}
\begin{theorem}[Hou--Feng--Leemans, Theorem 1.2]
For integers $n\ge10$, $s,t\ge2$ and $s+t\le n-1$, there is a
rank-three string C-group of order $2^n$ and type $\{2^s,2^t\}$.
\end{theorem}
\begin{proof}
Put $\delta=n-1-s-t\ge0$. The square group $H_{\delta+5}$ of
UA.Square has order $2^{\delta+5}$, pair orders $4,4,2$, and all
Boolean deletion maps. Attach marked dihedral quotients
$D_{2^{s+1}}\to D_8$ and $D_{2^{t+1}}\to D_8$ on its two adjacent
pairs. Each quotient is injective on its proper marked subgroups, and
the supports are the two pairs. UA.Edge gives a string C-group with
adjacent orders $2^s,2^t$ and order
\[
2^{\delta+5}\,2^{s-2}\,2^{t-2}=2^n.
\]
Distant commutation holds in every coordinate. The retained
group--polytope correspondence supplies the regular polyhedron and its
full automorphism group. This actually works for $n\ge5$ whenever the
same parameter inequalities hold; the displayed theorem keeps the
source's stated range.
\end{proof}

The remaining sections prove the complete classification in Theorem
1.3. For each $n\ge10$, the types $\{4,2^{n-3}\}$,
$\{4,2^{n-4}\}$ and $\{4,2^{n-5}\}$ have exactly two, two and four
ordered marked groups. Their complete presentations and all parameter
values are given below.
\section{Cyclic cores and the eight high-type families}\label{cyclic-cores-and-the-eight-high-type-families}

The supported constructions solve existence. Completeness at the high dihedral boundary needs a different local invariant, the cyclic core. Its bounded quotient is nevertheless recognized by the same square model, and its only finite seed is generated by the same UA.Nilpotent/UA.Ladder search as order 1024.

The retained \textbf{cyclic core theorem of Lucchini} says: if a cyclic proper subgroup \(M\) of a finite group \(H\) has core \(K\), then \(|M:K|<|H:M|\). Here \(K\) is the largest normal subgroup of \(H\) contained in \(M\). We use the form recorded in Isaacs, \emph{Finite Group Theory}, Theorem 2.20, also stated as L035 Proposition 2.9.

\textbf{Theorem H.Core (bounded cyclic quotient).} Let \((G;a,b,c)\) be a rank-three string C-group of order \(2^{t+s}\), type \(\{4,2^t\}\), with \(s\ge3\) and \(t\ge s-2\). Put \[
 r=bc,\quad m=2^t,\quad d_0=(ab)^2,\quad
 A=\langle r^{2^{s-2}}\rangle.
\] Then \(A\lhd G\), and \(G/A\) is a string C-group of order \(2^{2s-2}\) and type \(\{4,2^{s-2}\}\).

\textbf{Proof.} The even-word subgroup \(H=\langle ab,bc\rangle\) has index at most two. IP and UA.Phi give \(G/G'=C_2^3\), where its image has dimension two, so the index is exactly two. For \(M=\langle r\rangle\), the cyclic core theorem gives \(|M:\operatorname{Core}_H M|<2^{s-1}\), hence at most \(2^{s-2}\). Thus \(A\) is the subgroup of the cyclic core of its prescribed order, so is characteristic there and normal in \(H\). The involution \(b\) inverts \(r\); consequently \(A\lhd G\). It meets \(\langle a,b\rangle\) trivially by IP and lies in \(\langle b,c\rangle\). Their quotient parabolics still intersect in \(\langle bA\rangle\): equality of cosets gives membership in the original right parabolic, after multiplication by an element of \(A\). The left parabolic embeds and the right has its stated dihedral order. This proves IP, independence, type and order in the quotient. ∎

Since \(\operatorname{Aut}(A)\) is abelian, every commutator centralizes \(A\). In particular \([d_0,r^{2^{s-2}}]=1\). If \(A\ne1\), its unique involution \(z=r^{m/2}\) is central in \(G\). The universal identities \[
 r^a=d_0r,\quad d_0^a=d_0,\quad r^c=r^{-1},\quad d_0^c=d_0^r       \tag{A-C1}
\] follow from the involution, string and order-four relations.

Hereafter \(s\in\{3,4,5\}\), \(n=t+s\ge10\), so \(m\) is divisible by 32. The following eight models will realize every class. Let \(z\) be a distinguished central involution of the finite group \(D\):

{\def\LTcaptype{none} % do not increment counter
\begin{longtable}[]{@{}
  >{\raggedright\arraybackslash}p{(\linewidth - 4\tabcolsep) * \real{0.3333}}
  >{\raggedright\arraybackslash}p{(\linewidth - 4\tabcolsep) * \real{0.3333}}
  >{\raggedright\arraybackslash}p{(\linewidth - 4\tabcolsep) * \real{0.3333}}@{}}
\toprule\noalign{}
\begin{minipage}[b]{\linewidth}\raggedright
Branch
\end{minipage} & \begin{minipage}[b]{\linewidth}\raggedright
Group \(D\)
\end{minipage} & \begin{minipage}[b]{\linewidth}\raggedright
Right-conjugation automorphism \(T\), fixing \(z\)
\end{minipage} \\
\midrule\noalign{}
\endhead
\bottomrule\noalign{}
\endlastfoot
\(s=3\) & \(\langle d_0,z\rangle=C_2^2\) & \(T(d_0)=d_0z^\epsilon\) \\
\(s=4\) & \(\langle d_0,d_1,z\rangle=C_2^3\) & \(T(d_0)=d_1,\ T(d_1)=d_0z^\epsilon\) \\
\(s=5\) & Four involutions \(d_0,d_1,d_2,z\), \(z\) central, \([d_0,d_2]=1\), \([d_0,d_1]=[d_1,d_2]=z^\beta\) & \(T(d_0)=d_1,\ T(d_1)=d_2,\ T(d_2)=d_0d_1d_2z^\nu\) \\
\end{longtable}
}

Here \(\epsilon\in\{0,1\}\), or \((\beta,\nu)\in\{0,1\}^2\). For the last row an explicit group law on four binary coordinates is \[
 (u,v,w,e)(u',v',w',e')=
 (u+u',v+v',w+w',e+e'+\beta(vu'+wv')).                    \tag{A-C2}
\] Its bilinear final term proves associativity and gives all sixteen normal forms \(d_0^u d_1^v d_2^w z^e\). Thus the three \(D\)-orders are four, eight and sixteen. The displayed maps preserve their involution and commutator relations and have fourth power one. For example in the last row the two nontrivial commutator contributions cancel in \((d_0d_1d_2)^2\), and applying \(T\) four times to \(d_0\) returns \(d_0\); its orbit generates \(D\) together with \(z\).

Adjoin a cyclic \(r\) of order \(m\), with \(x^r=T(x)\), and take the central quotient \[
 N=(D\rtimes\langle r\rangle)/\langle zr^{-m/2}\rangle.       \tag{A-C3}
\] The identified subgroup has order two and meets both \(D\) and \(\langle r\rangle\) trivially, since \(T^{m/2}=1\). Hence \(D\) embeds, \(r\) retains order \(m\), and \(|N|=|D|m/2\).

Define automorphisms \(\mathcal A,\mathcal C\) by \[
 \mathcal A(r)=d_0r,\quad \mathcal A(d_0)=d_0,\qquad
 \mathcal C(r)=r^{-1},\quad
 \mathcal C(d_i)=T^{1-i}(d_0),\quad\mathcal C(z)=z.
\] In the first two rows \(\mathcal A\) fixes \(D\); in the third it fixes \(d_0,d_1,z\) and sends \(d_2\) to \(d_2z^\beta\). For an explicit verification, substituting the listed generator images into the defining involutions and commutators of \(D\) preserves each relation. On its generators the compatibility equalities are \[
 \mathcal A(Tx)=T(\mathcal A(x)^{d_0}),\qquad
 \mathcal C(Tx)=T^{-1}\mathcal C(x).
\] They express precisely preservation of \(x^r=T(x)\). Collection gives \((d_0r)^4=r^4\) in the first two rows and \(r^4z^\nu\) in the last. Since eight divides \(m/2\), this preserves \(r^m=1\) and \(r^{m/2}=z\). Finally \(\mathcal A^2=\mathcal C^2=1\) and \(\mathcal A\mathcal C=\mathcal C\mathcal A\) on \(r,d_0\), which generate \(N\). For instance both composites send \(r\) to \(r^{-1}d_0\). These checks prove the automorphism assertions for every admissible \(m\), not only for a computational sample.

Set \(G=N\rtimes\langle a,c\rangle\), where \(a,c\) are commuting involutions inducing these automorphisms, and put \(b=rc\). Then \(a,b,c\) are involutions, \([a,c]=1\), \(bc=r\), and \((ab)^2=d_0\). They generate \(G\), of order \(2^s m\), with pair orders \(4,m,2\). In \(G/N\), the images of \(a,b\) are independent whereas the right parabolic has image \(\langle cN\rangle\). Hence the only candidates for their intersection are \(1,b,d_0,d_0b\). The normal form (A-C3) gives \(d_0\notin\langle r\rangle\), leaving exactly \(\langle b\rangle\). The three generators are independent (already \(a\notin\langle b,c\rangle\), and the pair orders exclude the others), so this proves rank-three IP. Direct substitution yields \[
 \begin{array}{ll}
 s=3:&[d_0,c]=z^\epsilon,\\
 s=4:&[d_0,r^2]=z^\epsilon,\\
 s=5:&[a,r^2]^2=z^\beta,\qquad[a,r^4]=z^\nu.
 \end{array}                                                   \tag{A-C4}
\] They are invariant words of the ordered marking and distinguish the two, two and four models within their branches.

\subsection{Completeness and bare presentation recognition}\label{completeness-and-bare-presentation-recognition}

The eight source presentations have common relations \[
 a^2=b^2=c^2=(ac)^2=(ab)^4=r^m=1,\qquad r=bc,
\] and the following additional relations, with \(z=r^{m/2}\):

{\def\LTcaptype{none} % do not increment counter
\begin{longtable}[]{@{}
  >{\raggedright\arraybackslash}p{(\linewidth - 4\tabcolsep) * \real{0.3333}}
  >{\raggedright\arraybackslash}p{(\linewidth - 4\tabcolsep) * \real{0.3333}}
  >{\raggedright\arraybackslash}p{(\linewidth - 4\tabcolsep) * \real{0.3333}}@{}}
\toprule\noalign{}
\begin{minipage}[b]{\linewidth}\raggedright
Source name
\end{minipage} & \begin{minipage}[b]{\linewidth}\raggedright
\(m\)
\end{minipage} & \begin{minipage}[b]{\linewidth}\raggedright
Extra relations
\end{minipage} \\
\midrule\noalign{}
\endhead
\bottomrule\noalign{}
\endlastfoot
\(G_1,G_2\) & \(2^{n-3}\) & \([d_0,c]=z^\epsilon\), \(\epsilon=0,1\) respectively \\
\(G_3,G_4\) & \(2^{n-4}\) & \([d_0,r^2]=z^\epsilon\), \(\epsilon=0,1\) respectively \\
\(G_5,G_6,G_7,G_8\) & \(2^{n-5}\) & \([a,r^2]^2=z^\beta,\ [a,r^4]=z^\nu\), with \((\beta,\nu)=(0,0),(1,0),(0,1),(1,1)\) \\
\end{longtable}
}

\textbf{Theorem H.Classification (high-type completeness).} Every group in the three classes of L035 Theorem 1.3 is exactly one of the marked models constructed above; each displayed source presentation defines that model.

\textbf{Proof, first branch.} For \(s=3\), Theorem H.Core gives a quotient of type \(\{4,2\}\) and order sixteen, necessarily \(D_8\times C_2\). Thus \(h=[d_0,r]\in A=\langle r^2\rangle\). Since it commutes with \(r\), and \(d_0\) centralizes \(A\), \(h^2=[d_0,r^2]=1\). Hence \(h=1\) or \(z\), and (A-C1) makes this exactly \([d_0,c]=z^\epsilon\). The conjugates of \(d_0\) obey \(d_1=d_0z^\epsilon\); (A-C1) forces all the first model's actions and normal forms. The resulting surjection from that model onto the actual group is between groups of the same order, hence an isomorphism.

\textbf{Second branch.} For \(s=4\), the quotient of order 64 and type \(\{4,4\}\) is \(H_6\), by UA.Square\allowbreak Recognition. In that affine quotient \([d_0,r^2]=1\), so upstairs \(h=[d_0,r^2]\in\langle r^4\rangle\). Again \(h^2=[d_0,r^4]=1\), giving \(d_2=d_0z^\epsilon\), where \(d_i=d_0^{r^i}\). Conjugation by \(a\) fixes \(d_0,d_1\). Applying it to the equation for \(d_2\), with \(z\) central, gives \(d_2^a=(d_1^{d_0})^r=d_2\), so \([d_0,d_1]=1\). All second-model relations follow, and the same order comparison proves uniqueness.

\textbf{The bounded seed for the third branch.} The same universal-kernel theorem used for the order-1024 classification gives the following finite fact: every order-256 marked string C-group of type \(\{4,8\}\), with \(r=bc,d_0=(ab)^2\), satisfies exactly one row.

{\def\LTcaptype{none} % do not increment counter
\begin{longtable}[]{@{}llr@{}}
\toprule\noalign{}
Row & Relation pattern & Ordered marked classes \\
\midrule\noalign{}
\endhead
\bottomrule\noalign{}
\endlastfoot
I & \([a,r^2]^2=1,\ [a,r^4]=1\) & 1 \\
II & \([a,r^2]^2=r^4,\ [a,r^4]=1\) & 1 \\
III & \([d_0,r^2]=r^4\) & 1 \\
IV & \([(r^2)^a,r]=r^4\) & 1 \\
\end{longtable}
}

Apply UA.Nilpotent and UA.Ladder, as specialized in UA.Search with d=3 and k=3. The universal class-five group has order 2¹⁸ and third lower-central subgroup of order 2¹³. Its invariant index-two ladder has 7, 23 and 61 kernels. Testing all 64 subset pairs for each final candidate gives 37 ordered IP markings, four with the indicated orientation of type \{4,8\}. The census program in the shared computational appendix evaluates every displayed word on all four and verifies that each belongs to exactly one row, one marking per row. The complete program and every output check are printed in the shared computational appendix; UA.Search supplies the finite-search proof. Thus this bounded input is now conditional on exact GAP/NQ algorithms, with no abstract-group catalogue. No SmallGroups identifier or catalogue multiplicity enters these relation implications.

\textbf{Third-branch lifts.} For \(s=5\), Theorem H.Core gives \(A=\langle r^8\rangle\) and the table applies to \(G/A\). In row III we would have \([d_0,r^2]=r^e\), \(e\equiv4\pmod8\). It commutes with \(r\), so \([d_0,r^8]=r^{4e}=1\), impossible for \(m\ge32\), since \(v_2(4e)=4\). In row IV put \(x=(r^2)^a\). The relation would give \([x,r]=r^e\), \(e\equiv4\pmod8\), hence \(r^x=r^{1-e}\). The unit \(1-e\equiv5\pmod8\) has order \(m/4\) modulo \(m\): successive squaring gives \(v_2((1-e)^{2^j}-1)=j+2\). But \(x^{m/8}=(r^{m/4})^a\in A\), by normality, so it centralizes \(r\), contradicting that order. Thus III and IV cannot lift.

In I and II, \([a,r^4]\in A\), so \(a\) maps \(r^4\) to an odd power of itself and \(M=\langle r^4\rangle\lhd G\). Quotienting gives a marked \(\{4,4\}\) group of order 128, hence \(H_7=W/4\mathbb Z^2\). Let \(N=\langle d_0,r\rangle\). Equations (A-C1) make it normal, with index four: modulo \(N\), \(a,c\) commute, and their independence follows in the rank-three Frattini quotient, where \(N\) has the one-dimensional image spanned by \(r\). The cyclic normal subgroup \(M\) is central in \(N\), because \(d_0\) is a commutator and \(r\) centralizes its own powers.

In \(H_7\), the conjugates \(d_i\), \(0\le i\le3\), have linear part \(-I\) and translation parts \((0,0),\allowbreak(0,2),\allowbreak(2,2),\allowbreak(2,0)\), up to the choice of coordinate directions. They generate \(C_2^3\), with \(d_3=d_0d_1d_2\). Upstairs, every \([d_i,d_j]\) is in central \(M\) and has order at most two. Put \([d_0,d_1]=z^\beta\); conjugating by \(r\) gives \([d_1,d_2]=z^\beta\). Write \(d_3=d_0d_1d_2u\) with \(u\in M\). Conjugation invariance gives \[
 [d_1,d_2]=[d_2,d_3]=[d_2,d_0][d_2,d_1],
\] so \([d_0,d_2]=1\). Consequently \((d_0d_1d_2)^2=1\); because \(d_3\) is an involution, \(u^2=1\), giving \(u=z^\nu\). One further conjugation gives \(d_4=d_0\). These are precisely the final row of the table of the three model branches, and (A-C1) forces the actions of the commuting complement \(\langle a,c\rangle\). The model maps onto \(G\) and has its order \(32m\), so is isomorphic. This proves completeness of the eight models.

It remains to recognize the \textbf{bare presentations}, where the target order cannot yet be assumed. Let \(P\) be one of them and \(z=r^{m/2}\), initially of order at most two. For \(G_1,G_2\), if \(\epsilon=1\), \(a\) centralizes \(z=[d_0,c]\) since it centralizes \(d_0,c\); \(b,c\) centralize it since they invert \(r\). If \(\epsilon=0\), \(d_0\) is central; (A-C1) and the evenness of \(m/2\) then make \(z\) central too. Collection gives at most \(8m\) normal forms, attained by the first model.

For \(G_3,G_4\), \(d_2=d_0z^\epsilon\). If \(\epsilon=0\), applying \(a\) as in the second-branch argument gives \([d_0,d_1]=1\), after which the model relations and centrality of \(z\) follow. If \(\epsilon=1\), \([d_0,r^2]=z\) implies that conjugation by \(d_0\) inverts \(z\); as \(z^2=1\), it centralizes \(z\). Thus \(z\) commutes with \(r\) and all \(d_i\). Write \(j=[d_0,d_1]\). Applying \(a\) to the equation for \(d_2\), and using (A-C1), gives \(z^a=jz\). Since \(a\) fixes \(d_0,d_1\), it fixes \(j\), and \(a^2=1\) gives \(j^2=1\). Hence \(d_0,d_1\) centralize \(j\). Also \(j^r=[d_1,d_2]=[d_1,d_0z]=j^{-1}=j\), and \(c\) exchanges \(d_0,d_1\), so \(j\) is central. Collection yields \((d_0r)^4=j r^4\). Since \(m/8\) is even, \(z^a=(d_0r)^{m/2}=z\), so \(j=1\). The second model relations now give at most \(16m\) forms, all attained by that model.

For \(G_5\)--\(G_8\), put \(w=[a,r^2]\). If \(\beta=1\), conjugation by \(a\) inverts \(w\), and so fixes \(z=w^2\). If \(\nu=1\), it inverts \([a,r^4]=z\), again fixing \(z\). If both are zero, \(a\) centralizes \(r^4\) and therefore \(z\). Thus \(z\) is central in every case, and \([a,r^4]=z^\nu\) makes \(M=\langle r^4\rangle\) normal of order at most \(m/4\). Modulo \(M\), the presentation is \(W\) with the extra relation \([a,r^2]^2=1\). In the affine model \([a,r^2]\) is a doubled unit translation, so its square and conjugates kill precisely \(4\mathbb Z^2\). Thus the quotient has order 128, and \(|P|\le128(m/4)=32m\). The constructed model attains this bound and satisfies every relator, so the order comparison principle identifies \(P\) with it. This completes presentation recognition and the theorem. ∎

The preceding theorem proves the full marked interpretation of L035 Theorem 1.3; the listed word invariants distinguish all the declared families. It does not classify arbitrary two-groups of order \(2^n\), other types, or the source's unnumbered exceptional \(n<10\) discussion.

\section{Computation, attribution and formal scope}
The four shared files \texttt{phase78\_\allowbreak n009\_\allowbreak integration\_\allowbreak census.g},
\texttt{phase78\_\allowbreak n009\_\allowbreak integration\_\allowbreak finite.g},
\texttt{phase78\_\allowbreak n009\_\allowbreak integration\_\allowbreak data.g} and
\texttt{phase78\_\allowbreak n009\_\allowbreak integration\_\allowbreak run.py} are printed in full in the
shared appendix and distributed with this application. The first file
generates the order-256 seed independently of the four expected word
patterns. It checks every candidate and every subset intersection before
testing those words. Its exact output is 37 ordered markings, four of
type $\{4,8\}$, with one in each displayed row. GAP and NQ's exact
finite-group algorithms are retained computational machinery. The
complete files remain chargeable, including the portions reused by the
order-1024 application; no hypothetical smaller checker is counted.

The existence and classification results are due to Hou, Feng and
Leemans; the cyclic-core theorem is due to Lucchini. The common
supported construction and exhaustive kernel parameter space, together
with the explicit extension models and complete presentation bounds,
are the present derivation of those results. The calculations have been
reused from the reviewed proof bank with all hypotheses retained; no
historical priority for the derivation is asserted. The $n<10$ survey,
other types and a classification of arbitrary two-groups are outside
this paper's stated contract.

Existing Lean files establish specified kernel and intersection
interfaces in the common theory. They do not formalize the cyclic-core
instantiation, the eight models, their automorphism calculations, the
finite search, or the complete classification theorem. The written
proofs and exact computational obligation remain explicit; no new
source credit, question answer or length factor is claimed here.

\begin{thebibliography}{9}
\bibitem{UA} Human Compression Project, \emph{Supported kernels and
finite two-group quotients}, companion common paper, this volume.
\bibitem{HFL} D.-D. Hou, Y.-Q. Feng and D. Leemans,
\emph{Existence of regular 3-polytopes of order $2^n$},
arXiv:1802.08947v2, 22 January 2019; published as
\emph{2-groups behaving as automorphism groups of regular 3-polytopes},
Journal of Group Theory 22 (2019), 579--616.
\bibitem{Lucchini} A. Lucchini, \emph{On the order of transitive permutation groups with cyclic point-stabilizer}, Rendiconti Lincei: Matematica e Applicazioni, series 9, 9(4) (1998), 241--243. The cyclic-core form is also recorded in I. M. Isaacs, \emph{Finite Group Theory}, Theorem 2.20,
and in \cite[Proposition 2.9]{HFL}.
\bibitem{MS} P. McMullen and E. Schulte, \emph{Abstract Regular
Polytopes}, Cambridge University Press, 2002.
\end{thebibliography}

\end{document}
